\documentclass{article}
\usepackage[utf8]{inputenc}
\usepackage{amsmath,amssymb,amsfonts,amsthm,geometry}
\geometry{margin=1in}
\title{Comparison of the Right Shift Operator: Amann/Salamon vs. Axler}
\author{MathStudio Agentic Researcher}
\date{\today}
\begin{document}
\maketitle
\section{Introduction}
The Right Shift Operator is a fundamental example in functional analysis, particularly when studying the spectrum of non-self-adjoint operators on Hilbert spaces. This report compares the treatment of this operator in the works of Sheldon Axler and the functional analysis volume associated with the Amann/Escher series (Bühler and Salamon).

\section{Formal Definitions}

\subsection{Sheldon Axler (Measure, Integration \& Real Analysis)}
Axler defines the right shift $T: \ell^2 \to \ell^2$ on the Hilbert space of square-summable sequences as:
\[ T(a_1, a_2, a_3, \ldots) = (0, a_1, a_2, a_3, \ldots) \]
Source: Axler, \textit{Measure, Integration \& Real Analysis}, p. 314 (ID: 20865).

\subsection{Bühler and Salamon (Functional Analysis)}
In the context of the Amann series' transition to functional analysis, Bühler and Salamon define the operator $B: \ell^2 \to \ell^2$ similarly:
\[ Bx := (0, x_1, x_2, x_3, \ldots) \quad \text{for } x = (x_i)_{i \in \mathbb{N}} \in \ell^2 \]
Source: Bühler \& Salamon, \textit{Functional Analysis}, p. 223 (ID: 21237).

\section{Spectral Properties}

\subsection{The Spectrum $\sigma(T)$}
Both authors agree that the spectrum of the right shift operator is the closed unit disk in the complex plane:
\[ \sigma(T) = \{ \lambda \in \mathbb{C} : |\lambda| \le 1 \} \]

\subsection{Decomposition of the Spectrum}
A key distinction in the analysis is the decomposition into point, residual, and continuous spectra.

\subsubsection{Point Spectrum (Eigenvalues)}
Axler and Bühler/Salamon both state that the right shift has no eigenvalues:
\[ P\sigma(T) = \emptyset \]
This is because $T(a_1, a_2, \dots) = \lambda (a_1, a_2, \dots)$ implies $0 = \lambda a_1$, $\lambda a_2 = a_1$, etc., which for $\lambda \neq 0$ forces all $a_i = 0$.

\subsubsection{Residual and Continuous Spectrum}
Bühler and Salamon provide a more granular breakdown (p. 223):
\begin{itemize}
    \item \textbf{Residual Spectrum ($R\sigma$)}: The open unit disk, $\text{int}(\mathbb{D}) = \{ \lambda \in \mathbb{C} : |\lambda| < 1 \}$.
    \item \textbf{Continuous Spectrum ($C\sigma$)}: The unit circle, $S^1 = \{ \lambda \in \mathbb{C} : |\lambda| = 1 \}$.
\end{itemize}
Axler focuses on the total spectrum $\sigma(T)$ being the closed disk, while noting the lack of eigenvalues.

\section{Conclusion}
The definitions of the Right Shift Operator are identical across both sources. While Axler provides a concise summary of the total spectrum, Bühler and Salamon (representing the advanced trajectory of the Amann analysis series) provide the detailed decomposition into residual and continuous components, illustrating the operator's non-surjectivity and range density properties.

\end{document}